
\documentclass[11pt,a4paper]{article}
\usepackage[utf8]{inputenc}
\usepackage[T1]{fontenc}
\usepackage[english]{babel}
\usepackage[margin=2.3cm]{geometry}
\usepackage{booktabs}
\usepackage{amsmath}
\usepackage{lmodern}
\usepackage[protrusion=true,expansion=false]{microtype}
\usepackage[hidelinks]{hyperref}
\usepackage{parskip}
\usepackage{fancyhdr}
\pagestyle{fancy}
\fancyhf{}
\renewcommand{\headrulewidth}{0pt}
\renewcommand{\footrulewidth}{0.4pt}
\fancyfoot[L]{\footnotesize\href{https://github.com/gerontec/blastDiamond}%
{\texttt{github.com/gerontec/blastDiamond}}}
\fancyfoot[R]{\footnotesize\thepage}
\fancypagestyle{plain}{\fancyhf{}%
  \renewcommand{\headrulewidth}{0pt}\renewcommand{\footrulewidth}{0.4pt}%
  \fancyfoot[L]{\footnotesize\href{https://github.com/gerontec/blastDiamond}%
  {\texttt{github.com/gerontec/blastDiamond}}}%
  \fancyfoot[R]{\footnotesize\thepage}}
\usepackage{listings}
\lstset{basicstyle=\ttfamily\scriptsize, breaklines=true,
        breakatwhitespace=false, columns=fullflexible,
        frame=single, framesep=4pt, xleftmargin=0pt,
        postbreak=\mbox{\textcolor{gray}{$\hookrightarrow$}\space}}
\usepackage{xcolor}
\setlength{\emergencystretch}{3em}
\sloppy

\title{\textbf{Where the +1 frameshift peptides of BNT162b2 come from}\\[2mm]
\large A database check of Extended Data Table 3 of Mulroney et al.,
Nature 625, 189--194 (2024)}
\author{Analysis on a local mirror of NCBI nr, six DIAMOND part databases}
\date{@@datum@@}

\begin{document}
\maketitle

\section*{Summary}

This note reports an incidental finding. The starting question was narrow and
expected a negative answer: are the +1 frameshifted products described by
Mulroney et al.\footnote{Mulroney, T.E. \emph{et al.}
\emph{N1-methylpseudouridylation of mRNA causes +1 ribosomal frameshifting},
Nature \textbf{625}, 189--194 (2024),
\href{https://doi.org/10.1038/s41586-023-06800-3}{10.1038/s41586-023-06800-3}.}
present anywhere in NCBI nr? They are not, as the paper itself already
indicates. While answering it, something else turned up.

\begin{quote}
\textbf{Eight of the 123 peptides in the +1FS spike ELISpot pool do not come
from the spike reading frame at all.} Seven lie downstream of the spike stop
codon, in the human \emph{AES/TLE5}-derived element of the vaccine's 3' UTR,
and one lies upstream of the start codon in the 5' UTR. Their +1 translation
is therefore human sequence, which is why one of them matches a human GenBank
entry at 88\,\% identity over 50 residues.
\end{quote}

If this is right, a T-cell response measured against those pool members would
be directed at a peptide of host-derived sequence rather than at a foreign
product. Because responses were recorded per pool and not per peptide, the two
cannot be separated after the fact. We do not know whether this was intended,
and we may have missed a place where the paper addresses it.

Everything below is reproducible from public data; Appendix~B gives the exact
commands, accessions and two pitfalls that cost us time.

\section*{Provenance}

The analysis was carried out by Claude Opus 5, an AI assistant, working
against a locally mirrored NCBI nr. The finished document was then checked by
Claude Opus 5.5, which re-derived the numerical claims from the stored outputs
and corrected two of them: the count of genuinely novel peptides (115, not
118, once the eight UTR peptides are excluded) and the description of the 3'
UTR, where the mitochondrial element had been cited rather than verified and
the \emph{AES} element was mapped incompletely. Both corrections are
incorporated above.

It has not been reviewed by a domain expert and is not a manuscript. We are
sending it because the 3' UTR observation seemed worth a few minutes of the
authors' attention, not to make a claim about safety, which this work does not
address. If the point is already known or mistaken, please disregard it.

\section{What the paper reports}

\subsection{Mechanism}

N1-methylpseudouridine (m1$\Psi$) slows elongation and makes the ribosome slip
by one base at \emph{slippery sites}, defined in the paper as
\[
\text{m1}\Psi\;\text{m1}\Psi\;\text{m1}\Psi\;X,
\qquad X \in \{\text{m1}\Psi,\ \text{C}\},
\]
with $X$ the first nucleotide of the immediately downstream codon. Three such
sites exist in the Fluc reporter; site A precedes a +1 stop codon and
contributes nothing, while mutants U*187C and U*208C (sites B and C) strongly
reduce frameshifting and the double mutant abolishes it. Six identical sites
are distributed over the BNT162b2 spike CDS (Extended Data Fig.~6). The rate
is about 8\,\% in the reporter.

\subsection{Evidence in humans}

IFN$\gamma$ ELISpot on PBMC from 21 BNT162b2-vaccinated against 20
ChAdOx1-vaccinated donors, stimulated with the +1FS peptide pool. The response
is raised only after BNT162b2 ($P = 0.0233$, one-tailed Welch); in mice the
separation is clearer ($P = 0.0005$ against untreated).

\subsection{Sample selection and response rate}

Not randomised. The samples are a convenience set from four existing UK
collections: NIHR National BioResource (REC 17/EE/0025), PITCH as a substudy
of SIREN (IRAS 284460, ISRCTN11041050), a Liverpool protocol (16/NW/0170) and
the Oxford GI Biobank (16/YH/0247). For the animals the paper states it
plainly: \emph{``Mice were not randomly assigned to groups and experimenters
were not blinded''} ($n = 8$ per group).

Extended Data Fig.~2 encodes the per-donor outcome in the sample IDs, with a
threshold of $> 40$\,SFU/million. Counting them gives five responders
(155, 105, 51.7, 43 and 40 SFU/million) against 15 below threshold, i.e.\ a
response rate of about $5/21 \approx 24\,\%$. Two details are visible in that
figure: Donor-13 at exactly 40 does not clear its own ``$> 40$'' threshold yet
is flagged as a responder, and with five responders of 21 the one-tailed
$P = 0.0233$ rests on a narrow base.

\subsection{Degradation and inflammatory markers}

A full-text scan returns no occurrence of \texttt{degrad*},
\texttt{proteasom*}, \texttt{inflammat*}, \texttt{half-life} or
\texttt{misfold*}. There is no degradation kinetics, no CRP or IL-6 data and
no time course. Indirectly the ELISpot response does establish that the
products are degraded at all, since MHC-I presentation requires proteasomal
processing; it says nothing about how fast. A PubMed query
\texttt{frameshift AND (BNT162b2 OR N1-methylpseudouridine)} returns exactly
one record, the paper itself.

\section{Why nr holds no such product}

The \texttt{taxid} in nr describes the origin of the \emph{coding sequence},
not the cell that synthesised the protein. nr is a collection of submitted
records, and the taxon comes from the submitter's \texttt{/organism}
qualifier. A mistranslation product is coded by a synthetic, codon-optimised
construct; it is a translation error rather than a gene product, as the
authors establish by RNA-seq (Extended Data Table~2); and nobody submits one.
Only MassIVE \texttt{MSV000093074}, GEO \texttt{GSE223044} and figshare
\texttt{10.6084/m9.figshare.24271744} are deposited. The authors already
checked the database question themselves in Extended Data Fig.~4.

\section{The search}

\subsection{Query set}

@@class_summary@@

The positive control carries the design: firefly luciferase is in nr
(\texttt{P08659}, 550\,aa), so if its in-frame peptides are not recovered, the
search is misconfigured and a negative result means nothing.

The 123 peptides had to be transcribed from the raster image of Extended Data
Table~3. The pool is built from overlapping walks, which gives a
self-check: 41 transitions break (one per +1 stretch, each ending at a stop
codon) and the remaining 81 match letter for letter. One inconsistency
(\texttt{PPDSPLCTPGTGSGSAI} against \texttt{LCTPPGTGSGSAIAW}, one \texttt{P}
apart) is present in the published table and was left uncorrected.

\subsection{Parameters}

Six DIAMOND part databases covering about 1.18 billion sequences, built the
day of the run, NCBI state 23 Nov 2026. \texttt{-{}-dbsize} is set to the sum
over all parts so that E-values match a whole-nr search. Sensitivity is the
DIAMOND default, \texttt{-{}-evalue 1000},
\texttt{-{}-max-target-seqs 50}, 20 threads. The loose threshold is required
because peptides of 8--27\,aa score low; exactness is enforced afterwards,
counting only alignments with \texttt{pident} $= 100$ \emph{and}
\texttt{length} $=$ \texttt{qlen}.

As a sensitivity control, \texttt{p1\_mammalia} was repeated with
\texttt{-{}-ultra-sensitive -{}-masking 0 -{}-comp-based-stats 0}: 457\,s
instead of 18\,s, same result for the spike peptides.

\subsection{Wall clock}

@@runtime_table@@

Total @@total_runtime@@ for @@n_all@@ peptides against 1.18 billion sequences.
Run time follows part-database size, not hit count: DIAMOND reads the block
index linearly.

\subsection{Result}

@@n_raw@@ raw HSPs; after the exactness filter @@n_exact@@ rows remain, joined
against the local metadata tables:

\footnotesize
@@meta_table@@
\normalsize

Three observations. The positive control is recovered, so the negative result
for the spike peptides is meaningful. The MS-identified ``+1FS'' reporter
peptides are recovered too, which is by construction rather than a
contradiction: in Fluc+1FS the luciferase C terminus is deliberately placed in
the +1 frame, so a slip produces genuine luciferase sequence, which is why
enzyme activity returns. And none of the @@n_fs@@ spike peptides is found.

A side observation on database hygiene: those hits land not only on the
expected vector-derived record (one OID carrying 196 deflines, from
\texttt{pGL3}, \texttt{pGL4} and similar) but also on entries annotated as
genuine bacterial proteins -- \texttt{MBU6082087.1} \emph{Allobacillus
halotolerans}, \texttt{WP\_514138070.1} \emph{Streptococcus intermedius},
\texttt{WP\_125392797.1} \emph{Streptococcus gordonii}, all titled
``AMP-binding protein''. A 27\,aa peptide matching at 100\,\% is not chance;
this is vector sequence that entered those assemblies, and the title is
plausible because luciferase belongs to the acyl-CoA synthetase family.

A limit worth stating: DIAMOND is seed-and-extend and misses very short
queries. Control peptides \texttt{YGLNTNHR} (8\,aa), \texttt{KLPIIQK} and
\texttt{IIIMDSK} (7\,aa) were \emph{not} found although they are certainly in
nr. Among the spike peptides this affects at most three (two 8-mers and one
9-mer); the other 120 are $\geq 12$\,aa.

\section{Clustering against the human proteome}
\label{sec:cluster}

Exact search only answers identity. Whether a peptide is \emph{related} enough
to a human protein to group with it is a clustering question. Our daily nr
clustering compares new sequences only with each other; these peptides are in
no nr release, so that run would never touch them. The counterpart was
therefore restricted to \emph{Homo sapiens} (taxid 9606, 2\,285\,891
sequences), two levels at 90\,\% and 50\,\% identity.

One parameter had to change. The daily run uses \texttt{-{}-cov-mode 1}
(coverage of the target), under which a 15\,aa peptide can never join a
550\,aa protein because 80\,\% of the protein would have to be covered. With
\texttt{-{}-cov-mode 2} (coverage of the shorter sequence) a peptide can
become a member at all; otherwise every peptide would be a singleton by
construction.

@@cluster_section@@

\section{The closest thing nr has}
\label{sec:aehnlich}

Ranking the raw HSPs by bitscore answers the softer question. Of the 123 spike
peptides, \textbf{eight} have any hit at all, and the best five belong
together:

\begin{center}
\footnotesize
\begin{tabular}{llrrrl}
\toprule
peptide & coverage & identity & bits & E-value & subject \\
\midrule
\texttt{fs116} & 17/18 & 94.1\,\% & 42.0 & 4.5e--03 & \texttt{AAG23172.1} \emph{Homo sapiens} \\
\texttt{fs119} & 17/18 & 100\,\% & 41.6 & 6.4e--03 & \texttt{AAG23172.1} \emph{Homo sapiens} \\
\texttt{fs118} & 17/17 & 88.2\,\% & 39.3 & 4.8e--02 & \texttt{AAG23172.1} \emph{Homo sapiens} \\
\texttt{fs117} & 16/17 & 93.8\,\% & 37.4 & 2.7e--01 & \texttt{AAG23172.1} \emph{Homo sapiens} \\
\texttt{fs120} & 15/18 & 100\,\% & 35.4 & 1.5e+00 & \texttt{AAG23172.1} \emph{Homo sapiens} \\
\midrule
\texttt{fs082} & 15/15 & 66.7\,\% & 28.9 & 5.0e+02 & \texttt{KAF9461180.1} \emph{Collybia nuda} \\
\texttt{fs025} & 18/18 & 50.0\,\% & 28.5 & 8.1e+02 & \texttt{CEM26774.1} \emph{Chromera velia} \\
\texttt{fs006} & 15/15 & 66.7\,\% & 28.1 & 9.2e+02 & \texttt{KAF2594081.1} \emph{Brassica cretica} \\
\bottomrule
\end{tabular}
\end{center}

The bottom three are noise at E-values of 500 to 900. The top five are not.
\texttt{fs116} to \texttt{fs120} are five consecutive, overlapping pool
members; chained on their overlaps they give 50 contiguous residues:

\begin{center}
\texttt{SSWYCMHAMLAAPFPSWVPRVSPDLGSQVCSHLHLPHSPPLLVPDTSQAR}
\end{center}

Aligned without gaps against \texttt{AAG23172.1} (\emph{Homo sapiens},
103\,aa, gene \texttt{PP2150}, product simply ``unknown''):

\begin{lstlisting}
 +1                                       SSWYCMHAMLAAPFPSWVPRVSPD
                                           | |||||||||||||||| |||||
AAG   MGDQVPREQDLAKRSSPGPSPVLNMCNRQSAWATALSPWYCMHAMLAAPFPSWAPRVSPD

 +1   LGSQVCSHLHLPHSPPLLVPDTSQAR
        ||||||||||||||||||||||
AAG   PGSQVCSHLHLPHSPPLLVPDTSTPTWSSPIAHKRGGTRDELS
\end{lstlisting}

\textbf{44 of 50 residues identical, 88\,\%.} Against 1.18 billion sequences
that is not chance.

\subsection{Resolution: this is not spike}

\texttt{AAG23172.1} comes from \texttt{AF269289.1}, a 1035\,bp ``Homo sapiens
PP2150 mRNA'' submitted in May 2000, unpublished, product ``unknown'', with
the note \emph{``Method: conceptual translation supplied by author''} -- the
protein is itself only a guessed reading frame.

The nucleotides decide the matter. The 150\,bp of \texttt{AF269289.1} encoding
the matching stretch (nt 665--814) are \textbf{150/150 identical} to the human
\emph{TLE5} (formerly \emph{AES}) mRNA \texttt{NM\_198969.1} from position
1216, within its 3' untranslated region. The 3' UTR of BNT162b2 is built from
human parts: the 3' UTR of \emph{AES/TLE5} together with an element of
mitochondrial 12S rRNA (\emph{mtRNR1}). Both are verified directly in the
cassette below.

\subsection{Check against the construct}

We did not leave this at plausibility. GenBank holds the cassette as
\texttt{PP544446} (``Vector BNT162b2 clone Vacc\_Batch\_FP9359'', 7692\,bp,
deposited 2024); its spike CDS at nt 3578--7399 translates correctly to
\texttt{MFVFLVLLPLVS...KLHYT*}. Searching every peptide in the +1 frame
\emph{relative to the start codon} places them as follows:

\begin{center}
\begin{tabular}{lr}
\toprule
position in the cassette & peptides \\
\midrule
inside the spike CDS & 113 \\
\textbf{3' UTR, past the spike stop codon} & \textbf{7} \\
5' UTR, upstream of the start codon & 1 \\
\midrule
recovered exactly & 121 of 123 \\
\bottomrule
\end{tabular}
\end{center}

The seven downstream peptides are \texttt{fs116} to \texttt{fs122}, staggered
from +5 to +215\,nt past the spike stop; the upstream one is \texttt{fs001}.
Finally the 3' region itself. Mapping the 293\,nt following the stop codon
onto the two candidate sources resolves it completely:

\begin{center}
\begin{tabular}{llr}
\toprule
3' region & source & exact match \\
\midrule
nt 10--56 & \emph{TLE5}/\emph{AES} \texttt{NM\_198969.1} from 1221 & 47\,nt \\
nt 79--145 & \emph{TLE5}/\emph{AES} \texttt{NM\_198969.1} from 1290 & 67\,nt \\
nt 146--284 & mtRNR1 \texttt{NC\_012920.1} from 112 & 139\,nt \\
\bottomrule
\end{tabular}
\end{center}

Both human elements are therefore demonstrated directly in the cassette rather
than merely cited: the \emph{AES} element first, the mitochondrial 12S element
behind it. The gaps between the blocks correspond to the modifications in the
construct.

Placing the seven downstream peptides on this map assigns each to its source:
\texttt{fs116} (+5\,nt) through \texttt{fs120} (+101\,nt) fall in the
\emph{AES} part, which is why they match the \emph{AES}-derived entry
\texttt{AAG23172.1}; \texttt{fs121} (+125\,nt) straddles the junction, and
\texttt{fs122} (+215\,nt) lies in the mitochondrial element. The eighth
peptide, \texttt{fs001}, is upstream in the 5' UTR.

The chain is closed at every link: \texttt{fs116}--\texttt{fs122} sit past the
spike stop $\rightarrow$ the human \emph{AES/TLE5} element is there
$\rightarrow$ its +1 translation is 88\,\% identical to a human GenBank entry
whose coding nucleotides are in turn 150/150 from \emph{TLE5}.

\subsection{Transcription accuracy}

The construct check also validates the transcription: 121 of 123 peptides were
recovered letter for letter. \texttt{fs007}, \texttt{fs008}, \texttt{fs039}
and \texttt{fs123} each differ by a single residue (longest present fragment
11 to 15\,aa). For \texttt{fs039} this is the inconsistency already present in
the published table: \texttt{PGTGSGSAIAW} is in the construct, while
\texttt{LCTPPGTGSGSAIAW} is printed, one \texttt{P} too many. The other three
are most likely our own reading errors from the table image. None affects the
result: all four belong to the spike group, which has no hits either way.

\subsection{What follows}

\begin{itemize}
\item \textbf{Far from spike: necessarily so.} A frame shifted by one base
      shares no residue with the original protein; the same nucleotides
      regrouped give entirely different codons. That is arithmetic, not an
      observation.
\item \textbf{Novel: true for 115 of the 123 peptides.} They occur nowhere in
      1.18 billion sequences, not even distantly.
\item \textbf{Far from anything human: not throughout.} Eight peptides do not
      derive from the spike reading frame; seven of them translate the human
      3' UTR element. The ribosome is not inventing anything there, it is
      simply reading on.
\end{itemize}

\subsection{Amino-acid composition}

Across all 123 peptides (1983 residues), S, P, T, A and C account for
\textbf{64\,\%} of residues against roughly 30\,\% in ordinary proteomes.
Cysteine (7.3\,\%) and tryptophan (4.7\,\%) are about threefold enriched,
while aspartate and glutamate together fall to 2.4\,\% and tyrosine is
effectively absent (3 of 1983). This is low-complexity, proline- and
cysteine-rich sequence, the unavoidable consequence of reading a GC-rich
codon-optimised gene one base out of frame. Such sequence folds poorly and has
no natural counterpart, which makes the empty database result unsurprising in
hindsight.

\section{What this adds, and what it does not}

The added value over the publication is small and we would rather not oversell
it. The authors performed the database comparison themselves (Extended Data
Fig.~4). What is new here: a different tool against a much larger and current
reference (the whole of nr rather than two sequences compared pairwise); a
machine-checkable criterion instead of an alignment figure; a positive control,
which the paper does not carry; the clustering question; disclosed commands
and run times; and the 3' UTR observation. On that last point the paper's full
text contains no occurrence of \texttt{untranslated}, \texttt{UTR},
\texttt{AES}, \texttt{TLE5}, \texttt{mtRNR} or \texttt{12S}, so we believe
it is not discussed -- though the Methods do say the pool covers the +1 frame
of ``the BNT162b2 mRNA'', not of the spike CDS, so the coverage may well have
been deliberate.

What this does \emph{not} add: no new biology, no statement about risk or
tolerability, no measurement of the frameshifting rate, no clinical
interpretation. The expected negative result is confirmed. The contribution is
checkability rather than insight.

\section{Peptide list: +1FS spike (Extended Data Table 3)}

\footnotesize
@@peptides@@
\normalsize

\section{Appendix A: the queries as executed}

\subsection*{A.1 Invocation}

\begin{lstlisting}
fs_peptide_search.py fs_all.fa out_dir --mode fast
\end{lstlisting}

\subsection*{A.2 The DIAMOND query, once per part database}

\begin{lstlisting}
diamond blastp \
  -d parts/<part>.dmnd -q fs_all.fa -o <out_dir>/<part>.tsv \
  --outfmt 6 qseqid sseqid pident length qlen slen qstart qend \
             sstart send evalue bitscore staxids \
  --evalue 1000 --max-target-seqs 50 \
  --threads 20 --block-size 2 --index-chunks 1 \
  --dbsize @@dbsize@@
\end{lstlisting}

\texttt{-{}-dbsize} is the letter count over all six parts:

\begin{lstlisting}
SELECT SUM(letters) FROM blast_part;
\end{lstlisting}

The sensitivity control used the same fields but
\texttt{-{}-ultra-sensitive -{}-masking 0 -{}-comp-based-stats 0
-{}-evalue 10000}.

\subsection*{A.3 Timing and exactness filter}

\begin{lstlisting}
t0 = time.monotonic(); p = subprocess.run(cmd); dt = time.monotonic() - t0

if float(c[2]) == 100.0 and int(c[3]) == int(c[4]):   # pident, length == qlen
\end{lstlisting}

\subsection*{A.4 Metadata join}

\begin{lstlisting}
SELECT a.acc, a.taxid, s.len, s.title, t.name, t.rank
  FROM blast_acc  a
  JOIN blast_seq  s ON s.oid   = a.oid
  LEFT JOIN blast_taxon t ON t.taxid = a.taxid
 WHERE a.acc IN (%s, %s, ...);
\end{lstlisting}

One row per defline, so a summarised record carries many accessions and taxa
per OID; the \texttt{staxids} DIAMOND reports may differ from the taxid joined
here. For the luciferase record that OID carries 196 deflines.

\section{Appendix B: full reproduction}

Only \texttt{curl}, \texttt{python3}, \texttt{pdflatex}, DIAMOND, MMseqs2 and
access to nr are needed. Scripts are in
\texttt{github.com/gerontec/blastDiamond}, under
\texttt{examples/frameshift\_m1psi/}.

\subsection*{B.1 Fetch the sources}

\begin{lstlisting}
curl -o pmc.xml "https://eutils.ncbi.nlm.nih.gov/entrez/eutils/efetch.fcgi\
?db=pmc&id=10764286&rettype=xml"

B=https://media.springernature.com/full/springer-static/esm/\
art%3A10.1038%2Fs41586-023-06800-3/MediaObjects
curl -o Tab3.jpg "$B/41586_2023_6800_Tab3_ESM.jpg"   # 1398x2521
curl -o Tab1.jpg "$B/41586_2023_6800_Tab1_ESM.jpg"   # 2126x1065
\end{lstlisting}

Both tables are raster images only, so the peptides had to be read off them.
Columns were enlarged individually:

\begin{lstlisting}
convert Tab3.jpg -crop 380x830+618+1690 +repage -resize 300% c3.png
convert Tab3.jpg -crop 420x880+978+0    +repage -resize 300% c4a.png
convert Tab3.jpg -crop 420x880+978+860  +repage -resize 300% c4b.png
\end{lstlisting}

\subsection*{B.2 Validate the transcription first}

\begin{lstlisting}
check_overlap.py    # 123 peptides, 8-18 aa, no duplicates, 41 breaks
                    # (= 42 +1 stretches), 81 transitions match exactly
\end{lstlisting}

\subsection*{B.3 Build the query file}

\begin{lstlisting}
{ awk '{printf ">fs%03d_spike1FSpred\n%s\n",         NR, $0}' fs_peptides.txt
  awk '{printf ">ms%02d_fluc1FS_msproven\n%s\n",     NR, $0}' set2.txt
  awk '{printf ">ct%02d_fluc_inframe_control\n%s\n", NR, $0}' set3.txt
} > fs_all.fa       # 123 + 9 + 6 = 138
\end{lstlisting}

\subsection*{B.4 Search and clustering}

\begin{lstlisting}
fs_peptide_search.py  fs_all.fa out_dir    --mode fast
fs_peptide_search.py  fs_all.fa out_ultra  --mode ultra p1_mammalia
fs_peptide_cluster.py fs_all.fa cluster_out --taxid 9606 --threads 20
best_hits.py          out_dir
\end{lstlisting}

\subsection*{B.5 Resolve the striking hit}

\begin{lstlisting}
curl -o aag.gp "...efetch.fcgi?db=protein&id=AAG23172.1&rettype=gp"
#   -> 103 aa, gene PP2150, "unknown", coded_by AF269289.1:557..868
curl -o af269289.gb "...efetch.fcgi?db=nuccore&id=AF269289.1&rettype=gb"
align_aag.py   # 50 aa, best ungapped offset 36: 44/50 = 88 %
curl -o tle5.fna "...efetch.fcgi?db=nuccore&id=2462563547,39812018&rettype=fasta"
test_aes.py    # NM_198969.1: 150 of 150 nt identical from position 1216
\end{lstlisting}

Note: \texttt{NM\_005938} is \emph{not} AES but FOXO4. The correct accession
comes from \texttt{esearch.fcgi?db=nuccore\&term=TLE5[gene]+AND+"Homo
sapiens"[orgn]+AND+biomol\_mrna[prop]+AND+srcdb\_refseq[prop]}.

\subsection*{B.6 Check against the construct}

\begin{lstlisting}
curl -o bnt.gb "...efetch.fcgi?db=nuccore&id=PP544446&rettype=gb"
verify_utr.py  # spike CDS 3578..7399, +1 frame from the start codon
               # 113 peptides in the CDS, 7 in the 3' UTR, 1 in the 5' UTR
\end{lstlisting}

Two pitfalls that cost time when rebuilding this:

\begin{enumerate}
\item When reading the GenBank \texttt{ORIGIN} block, keep \emph{all} letters,
      not only \texttt{ACGTN}. Filtering out IUPAC ambiguity codes loses five
      bases here (7687 instead of 7692\,bp) and shifts the frame; the CDS then
      translates to \texttt{SCSWCCCLWCPA...} instead of
      \texttt{MFVFLVLLPLVS...}, which is easy to miss because it looks like a
      +1 product.
\item The relevant +1 frame is the one shifted against the \emph{start codon},
      not against the start of the vector. With the wrong reference point,
      none of the 123 peptides is found.
\end{enumerate}

\subsection*{B.7 Composition and report}

\begin{lstlisting}
composition.py     # S+P+T+A+C = 64 % of residues
fs_report_pdf.py out_dir fs_all.fa Frameshift.pdf cluster_out --lang en
\end{lstlisting}

\subsection*{B.8 Accessions used}

\begin{center}
\footnotesize
\begin{tabular}{lll}
\toprule
accession & what & used for \\
\midrule
PMC10764286 / PMID 38057663 & Mulroney et al., Nature 625 & source \\
PMID 38867495 & Boros et al., review & claimed follow-up \\
\texttt{PP544446} & BNT162b2 cassette, 7692\,bp & construct check \\
\texttt{AAG23172.1} / \texttt{AF269289.1} & PP2150, ``unknown'' & best hit \\
\texttt{NM\_198969.1} & TLE5 / AES mRNA & origin of the 3' UTR \\
\texttt{NC\_012920.1}:648--1601 & mtRNR1 / 12S & second UTR element \\
\texttt{MSV000093074}, \texttt{GSE223044} & MS and RNA-seq data & deposition \\
\texttt{P08659} & firefly luciferase & positive control \\
\bottomrule
\end{tabular}
\end{center}

\end{document}
